\documentclass{zhvt-classic}
\title{國劇身段譜}[國~劇~身~段~譜]
\author{齊如山}
\maker{甲辰年【龙】腊月十三酉时【日入】重製}

\begin{document}

\maketitle


\grid{none}{none}

\insertgraphic[width=0.65\linewidth,angle=90]{ti1}

\insertgraphic[width=0.65\linewidth,angle=90]{ti2}

\insertgraphic[width=0.65\linewidth,angle=90]{ti3}

\insertgraphic[width=0.65\linewidth,angle=90]{ti4}

\insertgraphic[width=0.65\linewidth,angle=90]{ti5}

\insertgraphic[width=0.65\linewidth,angle=90]{ti6}

\insertgraphic[width=0.65\linewidth,angle=90]{ti7}

\insertgraphic[width=0.65\linewidth,angle=90]{ti8}

\insertgraphic[width=0.65\linewidth,angle=90]{ti9}

\insertgraphic[width=0.65\linewidth,angle=90]{ti10}

\clearpage
\gridall

\tableofcontents
%  序

%  第一章 論剧來源於古之歌舞

%  第二章 論放劇與唐朝之舞有密切關係

%  第三章 論戲劇之身段

%    第一節【节】  袖谱

%    第二節【节】  手谱

%    第三節【节】  足谱

%    第四節【节】  腿谱

%    第五節【节】  腰谱

%  第四章論戲剧之身段

%    第一節【节】  鬍鬚谱

%    第二節【节】  翎子谱

\chapter*[序]{國劇身段譜句序}

\hfill 
佘研究戏中之身段、已二十佘年。从前对动作有不分明的地方、辄请问于梨园老辈俱蒙热心相告。凡有见闻、必记于小册、分类编成、日久积蓄颇多、恒欲撰为一书、以公同好。

去冬稍暇、即取小册、分类编成、又凭记忆所及、添入若干条。草就之后、即就正于梅君浣华、梅君又添出若干条。后又讲盆于徐君兰园、又为添出多条。于是将二君所说者、分别列入各类之中、书遂成。

但恐错谬太多。又求教于萧长华、叶春善、余叔岩、王瑶卿诸君、蒙其指正之处、亦复不少适【適】本会丛【叢】刊第二期出版、卽以此付印其中。名词及各角【脚、腳】身段之作法、当【當】仍有不切當之处【處】然以上诸君亦助余费若干心血矣。徐君兰园指示尤多、特书于此、以𰵧【誌】不忘、以后仍望高明、有以【敎】之也。

但有一节【節】余不能不郑重说明者：若人人都照这本书的写【寫】法去表演、纵然一丝不走、究竟千人一面、不但太科学化、有背美术【術】之原则、而且岂【豈qi】不太板滞、太机械、而毫无美趣了么翻过来再从一方面说、比方以抖袖一节【節】而论、可以算是板简单的身段了、但是古今诸位名脚抖的时候、是一人一个样、一出一个样、更可以说是一次一个样、而皆能美观。

果如此说、那规矩岂不没用了么【厯】却是我深信在当人手初学的时候、则必须照规矩来学、所谓不以規矩不能成方圆、以後則须神而明之、存乎其人、不可食古不化也。

民國二十一年春高陽齊如山識

\chapter*[第一章]{論剧來源於古之歌舞}[甲辰年腊月十六重制]
\begin{preface}
中国之戏剧、来源于古代之歌舞、早已为国人所公认、似乎不必再事翻缕矣。
\end{preface}

但此编为叙述戏剧典古舞递嬗之情形、故不惮烦琐、再中论之、中国自古公共庆【慶】贺或娱乐场所没有不歌舞的、有时或者不歌、然没有不舞的、此种记载、周朝以前、见于经传【傅】者已不少。汉魏六朝之间、歌舞风气较周朝又盛。试读其时之文字、便可知其大概。如观舞赋、舞赋、舞鹤赋、洛神赋、杯盘舞歌、白紵【zhu】舞歌、小垂手歌等等、难以尽【盡】述各篇中形容歌舞姿式之处、均穷【窮】形尽【盡】致、各极【極】其妙。固知其时、必定常有歌舞之事、入文人之目。故文人皆乐道之也。隋唐时歌舞之风尤盛、大致彼时得了许多外国跳舞的法子、又经皇帝提倡、朝野上下、都极【極】讲求。为【為】中国舞风极盛的时代。

故彼时歌舞的情形、见于诗章的更多、以上这些情形都与戏剧有直接的关系。为什么说古舞與现在的戏剧有直接的关系呢

一因古之娱乐场所、都有歌舞。今之戏剧完全是娱乐之事、且又完全是歌舞。所以说与古舞有直接的关系。

二因戏剧虽然是始白唐朝之梨园【園】但亦绝非是由唐朝骤然发生的。当然是由古代之舞、一朝一代的嬗变而来的。则与【與】古舞有直接的关系、是毫无【無】疑义【義】的。不过古时之歌舞简单、如今戏剧繁杂就是了。再说上边所说的历代文字、形容歌舞的情形、与现在戏中之规矩均极相合兹将书籍中形容歌舞的词句、略举数则、与见在戏中之粗织规矩、参证参证、便可知其有相常的关系了。

\begin{preface}
乐记﹁乐者、心之动也、声者、乐之象也、文采节奏、声之饰也。君子动其本、乐其象、然后治其节、是故先鼓以警戒。﹂云云。
\end{preface}

\begin{preface}
陈澔集说云﹁动其本、心之动也、心动而有声、声出而有文釆节奏、则乐饰矣。乐之将作、必先击鼓、以耸动众听、故日先鼓以警戒。舞之将作、必先三举足、以示其舞之方法。﹂云云。
\end{preface}

这一段文字、与戏中之情形、有极相似的地方。

比方戏中之起唱工、必系心内有所感触而起、或于自己思想事情的时候、或于听人述说事情的时候、心中有所感动、此时必定发出一种声音动作、作为 ﹁叫【呌】板﹂ 如云 ﹁呀哎﹂ 、或说 ﹁不好了﹂等等。或一抖袖、或一顿足、都可以作为叫【呌】板之表示。脚色一呌板、乐即起奏、这就是﹁乐者、心之动也。﹂的意思。
场面一听演员如何呌板、或喜怒、或哀乐等等、便知道该起什么样的唱工。此时锣鼓怎样的起法、琴笛等等怎样的奏法、哀有哀调、喜有喜调、这便是﹁声者、乐之象也。﹂ 唱工起后、琴笛怎样烘托唱者、休息的时候、音乐怎样的补垫、如大小过门等等。

再于歌唱之时、偶有身段、亦须添加一两下锣鼓、为的显着活动这便是﹁文采节奏、声之饰也。﹂

拉过门之前、必先起锣鼓、这便是﹁乐之将作、必先击鼓以怂众听﹂的意思。

演员出场之时、在门帘外必先有整衣转身等身段、后再往前行、这就是﹁舞之将作、必先三举足、以示其舞之方法﹂的意思。

\begin{preface}
  乐记日﹁诗言其志也、歌咏其言也、舞动其容也﹂
\end{preface}

戏中先是说白、说到心中有所感动之后、精神一震、便起唱工。

这个意思、〇只是说白觉得不解气了、非拉起嗓子唱一段、才觉养舒气痛快、这便是﹁歌咏其言也﹂。

起唱工之后、唱的词句是什么意思、必用手指点出来、唱到痛快淋漓的时候、必用手足全身来形容、他这些指点形容的举动、戏界都名之日身段、即是舞的原理、这就是﹁舞动其容也﹂。

\begin{preface}
  周礼郑注日﹁乐之声音节奏、未足以感动人而舞之发扬蹈【厲】为足以动人。﹂
\end{preface}

戏界最重作工、什么叫作工呢就是用手眼身法步、把唱歌句中的意思、都给形容出来、以便观客容易明瞭。所以成中于此地方、非常的注意、演员非精于此不可。比如演员如果作工好、就是嗓子不好、唱的差一点、观客也是欢迎的。他仍然可以出名。倘若该演员没有作工、专靠唱、那就非唱的极好不可。但他的享名、总不及有作工的演员欢迎的人较多。

本界的人、对于有作工的演员、尤为重视。因为有作工的演员、能够傅情、能够把本人的情绪、词句中的意思都给形容出来。观客看着容易提神、容易感动、所以大家特别的欢迎。比方贾洪林后半辈、就是全靠作工吃饭、这就是‘舞之发扬蹈厉【厲】为足以动人﹂的意思。

\begin{preface}
  诗序曰﹁咏歌之不足、不知手之舞之、足之蹈之、盖乐心内发、感物而动、不觉自运欢之至也.此舞之所由起也。﹂
\end{preface}

戏中常最初起唱的时候、往往不作身段、唱过一二句之后、虽作身段、亦只以手指示之。及至唱到酣畅淋漓的时候。便用股肱手足、全身的动作来形容。他在这个时候腔调板眼也快了、手足的动作、自然而然的也就快了。这便是﹁不知手之舞之、足之蹈之、乐心内发、感物而动、不觉自运欢之至也。﹂的意思。

\begin{preface}
  刘慷舞议日﹁舞之容、生于辞者也。﹂
\end{preface}

戏中于念白、或歌唱的时候、必须将制句中的意思、形容出来。如武家坡说：﹁全凭皓月当空﹂或唱﹁八月十五月光明﹂的时候、必伸两手作圆月式、以形容之。昆【崑】曲中对于此等地方、句句要形容、处处须认真、不必尽举。鄙人为梅兰芳所安身段、如奔月、散花、别姬等等、每一种舞式、均与词句关合、绝没有随便任意乱舞的时候。这就是﹁舞之容、生于辞者也﹂的意思

\begin{preface}
  傅毅舞赋曰﹁其始兴也、若俯若仰、若来若往、雍容惆怅、不可为象。﹂
\end{preface}

戏中演员、上场一出台帘【臺簾】先立住整整冠、糌绺须【鬚】抖抖袖、端端帯、欲行先不行、走一步又停一停、作种种姿式神情、这就是﹁其始兴也、若俯若仰、若来若往、雍容惆恨、不可为象。﹂的意思。由此更可知演员出场之后、时时刻刻、种种举动、皆系舞式也。

\begin{preface}
  舞赋又曰﹁其少进也、若翔若行、若竦若倾、兀动赴度、指颖应声。﹂
\end{preface}

戏中之快板及急急风的走法、或走同场、如穆柯寨穆桂英上高台之前、须前走一圆场、及各脚【腳jiao】﹁趟馬﹂之跑圆场等等步法、都是脚下要快、上身要稳、走的好的、真同飞【𮸽飝】一样、非常的美观。这就是﹁若翔﹂的意思。

戏中如趟马【趟馬】或出场、将要快走的时候、初一皋步、腳动的虽快、而地方却【卻】不大移动、又像前走、又不像前走、此种脚步在趟马【趟馬】时用的尤多。又如穆桂英在抬脸转身【臺臉𨍭身】后将上高台【臺】之前之几步、只举步不大移地方、也就是这种走法。这就是﹁若行﹂的意思。言其是欲行不行、似行非行也。合以上两种、也可以说又像飞【𮸽】又像走的让人捉揍不准。所以说是﹁若翔若行﹂这也有理。

再者演员每逢亮高像的时候、必定手往上伸、越高越好看、仿佛一块细长的石头【頭】立住的一样这便是﹁若竦﹂的意思。在亮矮像或卧于幕上、或斜身向前、或左侧右侧、彷佛将要倒下的一样、这便是﹁若倾﹂的意思。

言其是各种身段亮象、有时高有时矮、有时侧有时歪、种种变化、非常的好看。再者于演员作身段的时候、或稷立不动、或动之不已、〇之无论动与不动、或身段停止的时候、或立住初动的时候、都须有锣鼓交代、时时都要合拍、这就是﹁兀动【兀動】赴度﹂的意思。

演员于说白歌唱的时候、手是不断的动作、眼也不断的瞧看、但是手的动、眼的看、都要跟着腔调走。腔调快、眼就看的快、手也指的快。腔调慢、眼就看的慢、手也动的慢、腔偶止住、眼手也就跟着停住、并且是手指到什么地方、眼睛就看到什么地方、这是戏中一定的规矩。

所以说、手眼最要紧、这就是﹁指顾应声【指𮸹應聲】﹂的意思。

\begin{preface}
唐朝无名氏霓裳羽衣曲赋云﹁趋合规矩步中圆方。﹂
\end{preface}

戏中脚步、讲究方正、一步一步、都要踩的結实 转弯回身的时候、脚也有准尺寸、准方向不得一丝含混。可是跨腿转弯的时候又要圆活、不但脚步如此、一切身段、都须如此。

手一伸、足一拾、眼一看、身一转、无论何种身段、均要作到家。要作的磁实、所谓见梭见角、不可有一些含混。可是于动转的时候、处处都要圆、比方旦角用双手指的时候、无论左指右指、由手初动的时候、到手指定的时候、其间必规一圆圈式、就是一点小的曲折、也要圆活。于指定的时候、手却要指到确定的地方、头部眼神屑腰腿足都要与手随时动作、一毫不得或先或后、这叫作、圆处要圆、方处要方。这就是﹁趋合规矩步中圆方﹂的意思。

戏中演员刚出场的时候、必要在帘【簾】前停一停、于进场的时候、将身转的对了下场【塲】门的时候、未下之前、也要先停一停再走。此时虽然停住、但绝不能呆呆的立着。其精神彷佛有多们活动的意思。比方梅兰芳演宝莲灯【寳蓮燈】有王桂英出场的时候站在﹁九龙口﹂前边、等候胡琴拉过门完毕【畢】接唱一句慢板、大约总有两分钟工夫。在这个时候一动也不动、可要精神贯注、气足神完、有多少心事、都用眼神㳘露出来、仿佛有多们活动、有多少事情似的。令观客看着不但不呆板、且是非常的精种活潑。这就是﹁留而不滞﹂的意思。又戏中出场的时候、虽然用急急风的锣鼓、而演员也不能走的太快了、可是步步要有板。比方﹁辕门斩子杨六郎﹂出场、就是这个情形。又如穿靠的戏、于追下场的时候、虽用急急风、而演员也须脚步拿穏、脚根跺上劲、看着虽快、其实幷不快、而步步也都须把鼓点踩稳。又如断桥许仙走圆场的时候、走的脚步要稳、看着要快、虽然绕台一圏不过二三十步、可是看着仿佛跑了多远似的。又如陈德霖演探母回令在回令一场、太后出场的时候、锣鼓係用急急风、而太后穿氅衣高底鞋、怎么能够走的太快了呢。再说若跑的太快了、也不像太后的身分哪。德霖住这个时候、走的步步与锣鼓呼應、看着很快、其实并不快。以上这些情形、都是一﹁急而不促﹂的意思。古人云长袖善舞、现在衣服的袖子、都非常的长。袖子之外又加上一段水袖、更特别的长了。而戏中一切表情举动、又大致离不开袖子用袖子的动作、更不一样。或抖袖、或摔袖、或攘袖、或抓袖、或投袖、或擔袖、或扯袖、或咬袖、或挽袖、或举袖、或伸袖、或掉袖、或荡【盪】袖、或单或双、或反或正、一种有一种的姿式、各有不同。可是用袖子的时候、没有一处不跟腔调相合的、没有一处不跟锣鼓呼应的。无论怎样的动法、都要有板、这就是﹁纮无差袖﹂的意思。昆曲【崑曲】中无论那一出戏【齣戏、原指传奇中的一个段落、同杂剧中的﹁折﹂相近。今字作﹁出﹂指戏曲中的一个独立的段落或剧目】那一种腔调、均须与脚步呼应。如﹁夜奔山门﹂等等皆是。皮黄【皮黄又作﹁皮簧﹂是西皮和二黄的简称、它们是京剧的两大主要声腔、所以早年的京剧也被称为﹁皮黄﹂或﹁皮簧﹂戏】走边【走邊】【京剧﹁短打武戏﹂里的﹁走边﹂是一种重要的表演﹁程式﹂】的戏、唱吹腔的时候、也是处处与脚步呼应。就是皮黄中各种的腔调、虽然不能与步法呼应、也有许多地方、须要关合【關合】。

就如乌龙院、宋江出场唱四平调【由苏北花鼓演变而成的、在民间曾有﹁四拼调﹂的俗称】的时候也要步法间暇腔调安适、仍是彼此呼应的。比方此时若唱快板、那宋江就非快跑不可了。这也就是﹁声必应足﹂的意思。至于﹁万趋应矩同步中规﹂两句意义、已经在前面说明、不必再赘了。

\begin{preface}
沈郎霓裳羽表曲赋云﹁亦投袂而赴節﹂
\end{preface}

戏中袖子、最为重要、前已说明。总之袖子之举动、皆系音乐中之关節处。比方叫板时、常用抖袖、如起叫头【叫頭】时、则用举袖。如发狠时、则用抓袖。无可奈何时、则用荡袖。如决断时、则用投袖。如群英会打蓋后、周瑜下场之情形。不以为然的时候、则用摔袖。欢迎时候、则用搭袖。如此种种、难以盡述。总之袖子一举一动。都须介拍中節、都有锣鼓交代。内行所谓有﹁锣筋﹂是也。倘袖子动的时候不在板眼上、不在锣鼓点上、就是走板、那就显着松懈多了。这就是﹁亦投袂而赴節﹂的意思。古来书籍中、关于歌咏、或议论歌舞文字、类似这种的诗句、不知有多少、抄也抄不清、䤸也䤸不尽。以上不过每时代略举一二条、藉他来参证参证就是了。按古人这些词句、都是形容各该时代之歌舞的文字、并不是再与现在戏剧作的文字。
可是拿他的形容词来与戏剧的身段一参证、是非常的吻合。这些词句、就彷佛尃为现在戏剧之身段作的一样。这足见戏剧之身段、与古来之舞、是有直接关系的了。戏中的身段步法表情、以及种种动作、都是由古舞嬗变而来的、是毫无疑义的了。以上所引的、都是唐朝以前的文字、到宋朝以后、凡是歌咏、或记载歌舞的文字、差不多是具体的、与现在戏剧一样也就不必引证了。

\chapter*{第二章}{論放劇與唐朝之舞有密切關係}   

现在戏剧的身段、固然是由古代之舞嬗变而来。但是由唐朝起、才有了具体的规模。

\begin{preface}

涵虚子论剧云﹁戏剧至隋朝始盛、隋谓之康衢戏、唐谓之梨园乐、宋谓之华林戏、元谓之升平乐﹂云云。

\end{preface}

\begin{preface}

陶九成论曲云﹁唐有传奇、宋有戏曲、金有院本雑剧、而元因之。然院制杂剧、厘而为二矣﹂云云。

\end{preface}

以上二说、一论戏剧、一论词曲、意思大致相同。涵虚子虽然说是戏剧至隋朝始盛、然隋朝为期甚短、不见得有完备的组织、当然是到了唐朝、经有大力量的唐明皇一提倡、才有了完备的规定。以后由宋而元而明清、虽一朝与一朝的名目不同、但是没有什么极大的变更一直传流到了现在、各朝之命名、虽然屡经更改。【清明管戏剧的衙门亦名曰昇平*】而戏界至今、仍日【梨园行】就是戏中的各种身段步法、大致也来源于唐朝之舞。按古人歌舞曲牌的名子、存留者尚不少、而怎样的舞法、各种记载中、则毫未提及、无从参考。惟有宋朝周密的癸辛杂识里头、有一段记载说是﹁予尝得故都德寿宫舞谱二大帙、其中皆新制曲、多妃嫔诸阁、分所进者。所谓讲者、其间在所谓：

~左右垂手~雙拂~抱肘~合蝉~小𨍭~虚影~横影~稱裹 

~大小转攛~盤转~叉腰~捧心~叉手~打场~挽手~鼓兒

~打鸳鸯场~分颈~回头~海眼~收尾~豁头~舒手~布过

~鲍老惙~ 对窠~ 方腾~ 齐收~ 舞头~ 舞尾~ 呈手~ 关宝


~掉袖儿~~ 拂~~ 躜~~ 绰~~ 觑~~ 掇~~ 蹬~~ 焌


~五花儿~~ 踢~~ 搕~~ 剌~~ 𭣇~~ 繫~~ 搠~~ 捽


~雁翅儿~~ 靠~~ 挨~~ 拽~~ 捺~~ 闪~~ 缠~~ 提


~龟背儿~~ 踏~~ 儹~~ 木~~ 摺~~ 促~~ 当~~ 前


~勤蹄儿~~ 摆~~ 磨~~ 捧~~ 抛~~ 奔~~ 抬~~ 撅

以上唐代之舞谱也。其中所云之左右垂手、大小𨍭攛、打鸳鸯场等九种名目。乃是舞牌之名。至雙拂、盤𨍭、分颈、对窠等等六十三种名目、乃各种姿式之名词也。只看他将各种舞的姿式、又特定名目、便可知彼时歌舞之盛行、及衆【众】人研究之状况。当年此种舞谱、一定很多、可惜存留到今者、只有一二页、可谓吉光片羽矣。但是我们拿这几十个名词、与现在戏中之身段来参证、就觉着已经有许多相合的地方。是可知戏剧之身段、实来源于唐代之舞也。现在把谱中姿式的名词、与戏中各种身段姿式来参证参证：

~~左右垂手~~ 这个舞牌中姿式的名词、于现在戏中旦脚出场的各种身段、大致相似且读古来各种记载和歌咏的文字、大致都说这个舞牌子是女子所舞的；则旦脚出场之身段、来源于此也未可知。因旦脚所去的都是女子也。

兹参证如下比方：

~~双拂~~ 即戏中之双抖袖、因拂字之义、舆抖字大致相同也。

~~抱肘~~ 即戏中之合抱、俗名抱夹、乃双手交于胸际。

~~合蝉~~ 即戏中之双背袖即双贫手、刚袖往后一背、如同蝉之合翼。

~~小转~~\begin{preface}
即戏中之半转身、即往后看。按宋朝朱熹诗傅释﹂辗转反侧四字之意、曰辗者转之、转者辗之周、反者转之过、侧者转之留﹂舞谱亦用此字样、日转初势、转半势、转周势、转过势、转留势、伏覩势、仰瞻势、回头势。摢此则小转之意、当然要外乎回头、和半转身的情形。  
\end{preface}

~~虚影~~\begin{preface}
即戏中之掍、或一掍、或两掍、即然曰虚影、当然不是呆板的姿式、所以说他是戏中之掍。
\end{preface}

横影
~~即戏中之又舒袖、即两手旁伸。從前旦脚此種身段颇多、梆子班中、花旦尤多。又手往旁一伸、则横*于高、与横影二字的意思𣑕合。

稱裹
~~即戏中之左右理装、俗名左右看、当看左旁的时候、则將右袖举起、用左袖作担鞋现形式。看右旁的时候反是、所谓稱裹者、大致係看看装裹的相稱、不相稱的意思。不则裹字在姿式名词中、用之殊觉生硬。

掉䄂儿
~~这个舞的牌子、按古人文字中考之、也是女子所舞的、且与旦脚出场之各種身段、也大致相似。比方：

拂
~~~即戏中之單【单】抖袖。

躦
~~~即戏中之蹲、類篇云﹁𧿙、聚足、或作蹲。﹂又云﹁踏也。﹂戏中每合抱、必蹲身子、證以類篇所谓聚足又踏也等语、或亦係旦脚之踏步。即左足斜置于右足之后、而两腿聚靠、或右足斜置于左足之后、旦脚亮像必须如此。

綽
~~~~即戏中之欲前走后退一二步之身段、所谓绰約也。

覷
~~~~即旦脚之用手按鬓角、而往前看之姿式。旦角出场立定时、必有这种身段。

  掇

即戏中之提鞋式、或蹲身大抖袖式、此为旦脚出场必有的身段。用右手在右腿后提左鞋、用左手在左腿后提右鞋、花旦则用手掇其鞋后根、青衣则略以见意而已。

  蹬

即戏中之超步、所谓𨃴蹬也。

  焌

疑即戏中之双举袖前行式、旦脚行路、必打腰裙、在匆或着急的时候、必将袖双举。比方牧羊圈讨饭一场、青衣与老旦出场时、青衣之身段、便用此式。所谓燃烧火*上之炎状也。又周禮注﹁杜子春曰：焌與俊通。﹂则亦是正面与人看之义。此二釋義、与身段都算相合。

大小转撺

这个舞牌、在记载中不多见、揣摩他舞的情形、似一种群舞。但是各人的身段、都是相同的。且其名词之情形、亦与日脚之身段、有相似的地方。但此与前左右垂手、掉袖儿两种有不同之处。盖前两种之姿式、差不多与旦脚出场之情形、完全相同、可以说旦脚出场之姿式、完全是采用的那两种舞、至于此种舞牌、虽与旦脚之身段、有相合之处；但系散见于各戏之中、不似前两种之整个的采用也。现在也把他来参证参证：

  盘转

疑即戏中之绕圆场、现在仍名日走圆场、或跑圆场。如穆柯寨穆桂英登高台之前、及射雁后进场之前、都有此种身段。且绕圆场三字的意思、与盘转二字也相同。

  叉腰

即戏中之叉腰、双手叉于腰间、手背朝裹、肘往前伸、此种身段、旦脚用之、非常美观。

  捧心

即双手重叠于胸前、现在仍名日捧心、但亦旦脚用之。

  叉手

此名词颇难解释、疑系多人互相叉、手戏中不多见。如以个人叉手解释、则係双手互叉、手心向下之状。旦脚往往有此身段、花旦尤多、但此不过在疑似之间。

  打场 

疑即戏中之旋转、因*家打场、皆係牵引碜碡旋转也、河北童谣、有打场谣曰﹁咧咧咧咧场、打喽𢾶君打高梁。﹂等语。小儿每逢念此、身子必旋转、而每旋转、亦必歌也。此谣行之颇广、来源当亦甚远。故疑打场二字、即旋转之义。戏中旋转之处甚多、亦有两人互转之时。二人各以左膊相跨而转、如金山寺青蛇白蛇、及牧羊圈青衣老旦讨饭一场、都有这种身段。

  攙【搀】手 

疑即戏中二人两手互搀、如金山寺青蛇白蛇、及奇双会赵宠合李桂枝等等、都有此身段。惟此种身段、有相哭时方用之。

  鼓儿

阙疑

  勤蹄儿

这个舞牌、虽然以蹄字为名、但谱中名词、除奔字身外、大致都是用手的名词、所以鄙人疑惑此种舞系用袖子之舞、因谱中名词与现在戏中用袖子之之词大致相同也。

比方：

  擺【摆】
  

即戏中之摆袖、戏中的名词、用袖先往裹一转、再往外摆者名抖袖。一直往外一摆者、日名摆袖。

  磨

即戏中之磨袖、似担靴的样而确非担靴、比方以左袖平举而拍右腿、即以右袖绕转于腿之旁。如张飞及火判、常有这种身段。

  捧 

即戏中之捧袖、如双袖先往裹一折、作拱手状、此即名日捧袖。

  抛

即戏中之抛袖。抛袖之形势、不止一种、如双袖同时往前一挥、或两袖错棕挥之。总之戏中之抛袖、投袖、挥袖、掷袖等等、都是抛的性质、此说参看第三章。

  奔 

此字颇难解释、或係往前抛袖取奔放之义、因若以字面解释则只有快跑一解。在一小小舞场之中、也似乎用不着快跑的身段。

  擡

即戏中之撩袖、于惊讶、恐惧、大哭时、往往用此身段。其式亦不一、详见第三章。

  撅

即戏中拈袖之意、戏中有抓袖、拈袖、折袖、攘袖等、皆系撅袖的意思。详见第三章。

  龟背儿

这个舞牌、在记载中、亦不多见、不知是什么意思、按唐朝段安节乐府雜录中、有：字舞、花舞等名目系用舞的人、排成一个字、或排成一朵花的意思。如此则此舞牌、或者是排成一个龟背的形式、也未可知。但看此谱中的名词、亦系多数人的舞法、好像与现在戏中堆花、堆鬼、摆莲花灯、摆七巧阃等等的舞法、大致相同。

踏

疑即戏中之踏步。右足迈于左腿之左、或左足迈于右腿之右、皆名曰踏步。此系生净等男子的身段。旦脚踏步尤其多遇亮像时、必须踏步、就是将左足聚立于右足之后、或右足坚立于左足之后、都是踏步。再者戏中多数人舞时、往往此人跪于地上、又一人用足踏于跪者之腿上、此在群舞中乃极常用之身段。或与此踏字有直接关系也未可知。

儹

正韵云﹁儹者、聚也。﹂疑即戏中之攒字、因儹攒二字、音义都相同也。如今成中、只要大家在一处凑聚、便名曰攒。武行起打、一人挡四人的时候、就叫结攒；俗名又叫结个攒、则攒字完全成一名词矣。与此儹字必有相当的关系也。

木

疑系直立如木的意思。戏中多数人的身段、往往彼此对立、亦名日亮像。

摺【折】

即戏中之翻。如数人舞、一行直走的时候、忽然自第一人起、往回折走、即名日翻；或日回翻、即系折叠的意思。按广韵曰﹁摺亚也。﹂或系叠架之意。戏中也有这样的舞法、如堆灯、堆花等等、往往将切末叠于一处、就是一排比一排高的意思。总之以上两种解释法、一为横叠、一为竖叠、大致都与戏中舞式相合。

促

集韻曰﹁迫也；近也；密也﹂如此解法便是往一处凑集的意思。按戏中群舞、如摆七巧图、或四人武舞等等的时候。有时众人提四门斗、凑在一处、名曰合拢；亦曰凑。若分开、则名曰拉开。按合拢之名词、舆促字的意思颇相合。

当

此字当然是舞者立定、以迎来者的意思、与戏中拉开时之身段颇相合。

當

此字当然是舞者立定、以迎来者的意思、与戏中拉开时之身段颇相合。盖拉开时、必各拉山膀、又武舞。亮像之形势与此亦合。

前

此当然是一齐向前的意思、或者就是戏中之堆花、一排一排的依次向前走的样式。

雁翅儿

舞牌名曰雁翅、想来是所舞的姿式、彷彿大雁的翅膀一样的意思；各种姿式、即像翅膀、则舞者大致不是一个人、其数必在一人以上。且看其舞谱名词、也像两个人、或幾个人合舞的意思、于戏中两人同作的身段、相合的地方很多。而转、姿式亦极美观。或彼此用肩相靠。戏中此种身段亦颇多。

挨

说文﹁挨、击背也。﹂又正字通﹁凡物相近、谓之挨。﹂扬子方言谓﹁强进曰挨。﹂按说文是打的意思、正字通与扬子方言是挤的意思。戏中二人对舞时、用手相击、名曰搕。如金山寺、青蛇白蛇对唱时、恒彼此以剑对搕等等形势合意、都与正字通的解释法相同。愚常疑此种身段之不甚合道理、二人既非互打、又非练习、为什么要彼此相击呢？其来源大致与此有关。又戏中旦脚、用双手叉腰、横行俏步、二人对行、行至一处、两肩相摩、即名曰挨、如佳期红娘唱﹁待红娘轻轻悄悄把门挨﹂时、即系双手叉腰、俏步横行、想这就是形容挨字的意思。这种情形、与正字通及方言之解释法、亦大致相合。

拽

即戏中二人、此人用右手拉彼人之左手、同作卧雲式之身段。因此种身段、实是彼此相拽、且所作身段之姿式、亦如鸟之展翅、与舞牌之名亦相合。

捺

捺者、按也。疑即放中二人之亮高矮像。此人蹲下、彼人立于后边、以手按前人之肩。群舞之时、此种身段尤其多。

闪

即戏中之闪开。比方二人同作身段、走至相近之处、必要往外一躲、又名曰往外一闪。

缠

疑即戏中之二人拉手、同作鹆子翻身、因此即系缠绕的意思。再如二人用枪打仗时、两枪头互搅、即名曰缠、亦名曰绕。

提

戏中蹲身之后、身子起的时候、腿亦随之提起、即名日提。或此人按彼人之肩、身子往上一竄【窜】亦名曰提翻筋斗、有一种名单提。以上三种或与此提字有点关系。

鲍老惙

鲍老惙三字之义意、不大明了、但其中之名词、则完全是多数人合舞之意、与戏中四人合作之身段情形、有相同之点。

对窠

即戏中数人、合凑而各作同鲜身段的意思。

方腾

比名词在戏中、亦颇普通、与四门斗亦大致相同。

齐收

疑即戏中了舞时、各囗本位的意思、戏中即名曰囗本位。按囗本位三字、与齐收二字、亦大致相同。

舞头

阙疑。

舞尾

阙疑。 

呈手

疑即戏中群舞时、各作供手之式。

关宝

此二字、殊无索解。

五花儿

这一种舞牌、一定是各持器械、与戏中之舞如对枪、对刀等场子之姿式、相合。

踢

或即係戏中二人相踢、或一踢一拍。

榼

戏中用刀枪、彼此相撀、名曰搕。

刺

戏中用枪直刺方之人、即名曰刺。

𭣇

戏中甲用枪刀、压住乙之枪刀、乙必用手极力𭣇之、此即名曰𭣇。再一人持枪、以持枪之手背、撀左手手心、亦名曰𭣇。此种𭣇法、皆用三次。

繫

疑即戏中对枪时、彼此用枪头一绕、两绕、三绕。此种打法、含有彼此枪头、繫缠在一处、不得离开的意思。按这种情形、与繫字或有相当的𨶹係。

搠

即戏中扻或砍的意思。

捽

即戏中揪的意思、如甲揪乙之领、乙即翻一筋斗、即为捽字之原義。

打鸳鸯场

此种舞谱之名词、以皆在鸳鸯身着的意、戏中似无此种身段、不必强为拉扯。

按以上所说的情形、惟免有望文生训的地方、读者必定笑鄙人强为拉扯附会。然而一个人只有两手两足、古今動作、当亦无大差别、如此解法、虽然不敢说一定对、但是相去也不到其遠。由此更可能證明有机会在戏中之身段、實来源古时之舞、则毫无疑义矣。

\end{document}